\begin{scanpage}{0008}{iv}
\begin{center}FOREWORD\end{center}
First of all, one should explain the term ``physical theory of diffraction''. In order to do this, let us discuss briefly the historical development of diffraction theory.

If one investigates, for example, the incidence of a plane electromagnetic wave on a body which conducts well, all the dimensions of which are large in comparison with the wavelength, then the simplest solution of this problem may be obtained by means of \uline{geometric optics}. It is known that in a number of cases one must add to geometric optics the laws of \uline{physical optics} which are connected with the names of Huygens, Fresnel, Kirchhoff and Cotler. Physical optics uses, together with the field equations, the assumption that in the vicinity of a reflecting body geometric optics is valid.

At the start of the Twentieth Century, a new division of mathematical physics appeared --- the \uline{mathematical theory of diffraction}. Using it, rigorous solutions to the problem of diffraction by a wedge, sphere, and infinite cylinder were obtained. Subsequently, other rigorous solutions were added; however, the total number of solutions was relatively small. For sufficiently short waves (in comparison with the dimensions of the body or other characteristic distances) these solutions, as a rule, are ineffective. Here the direct numerical methods also are unsuitable.

Hence, an interest arose in approximation (asymptotic) methods which would allow one to investigate the diffraction of sufficiently short waves by various bodies, and would lead to more precise and reliable quantitative results than does geometric or physical optics. Obviously, these methods must in some way be considered the most important results extracted from the mathematical theory of diffraction.
\end{scanpage}
